% Paired with coordinate-frame transformations.
\clearpage
\section*{Verification case: a gradient depends on the chosen ruler}
\noindent Coordinate derivatives are not automatically vector components.
Let a world-space point be \(p=Aq\), with
\[
A=\begin{pmatrix}2&1\\0&3\end{pmatrix},
\qquad f(p_x,p_y)=p_x+2p_y.
\]
The differential in world coordinates is the linear functional
\(df_p(\delta p)=\begin{pmatrix}1&2\end{pmatrix}\delta p\).
With \(p=Aq\), a local coordinate perturbation obeys
\(\delta p=A\,\delta q\), and therefore
\[
df_q(\delta q)=\begin{pmatrix}1&2\end{pmatrix}A\,\delta q
=\begin{pmatrix}2&7\end{pmatrix}\delta q.
\]
The numbers \((2,7)\) are the \emph{covector} coefficients of the local
differential, not the coordinate components of the same geometric
steepest-ascent arrow. The local basis is slanted and unequally scaled,
so local-coordinate dot products do not measure world-space lengths.

The pullback Euclidean metric is
\[
G=A^TA=\begin{pmatrix}4&2\\2&10\end{pmatrix},
\qquad
\lVert\delta p\rVert^2=\delta q^TG\,\delta q.
\]
A metric-gradient vector \(g_q\) is defined by
\(df_q(\delta q)=g_q^TG\,\delta q\) for \emph{every}
\(\delta q\). Hence
\[
Gg_q=\begin{pmatrix}2\\7\end{pmatrix},\qquad
g_q=\begin{pmatrix}1/6\\2/3\end{pmatrix}.
\]
Multiplying by \(A\) gives the world vector \(Ag_q=(1,2)^T\),
which matches the standard Euclidean gradient of \(f\).

\paragraph{Why this matters.}
An optimization routine that takes the coordinate derivatives \((2,7)\)
as an ordinary direction in the slanted local coordinates changes its
geometric steepest-ascent meaning. In an orthonormal frame \(G=I\) and
this distinction is invisible; that convenient special case must not
be mistaken for the general rule. Surface gradients, covariance
ellipsoids, constrained optimization and normal transformations all
rely on the same covector--metric distinction.

\paragraph{Reconstruction exercise.}
Implement \(A\), \(G\), and the two linear maps in C++23 and Python.
For five non-axis-aligned perturbations \(\delta q\), check both
\(df_q(\delta q)=\begin{pmatrix}2&7\end{pmatrix}\delta q\)
and \(df_q(\delta q)=g_q^TG\delta q\).
Then replace \(A\) with nearly parallel columns, calculate
\(\kappa_2(G)\), and explain why the coefficients become sensitive.
\emph{Do not} claim that a small residual alone proves accurate recovery.
